\chapter{Implementación}
\label{ch:implementacion}

\section{Estructura del código y componentes}

El paquete \texttt{seshat} se organiza en cinco módulos (figura~\ref{fig:package-tree}):
\texttt{core/} contiene los modelos, configuración y utilidades; \texttt{infra/} agrupa
los adaptadores de sistemas externos; \texttt{app/} es la capa de aplicación;
\texttt{eval/} contiene los arneses de evaluación, que no forman parte del sistema en
producción; y \texttt{cli/} expone la línea de comandos.

\newcommand{\dtc}[1]{\hfill\makebox[10cm][l]{\textcolor{gray}{\texttt{\small\# #1}}}}

\begin{figure}[H]
\dirtree{%
  .1 seshat/.
  .2 app/.
  .3 agents/.
  .4 identification/.
  .4 resolution/.
  .3 pipeline/.
  .4 extraction/ \dtc{dos pasadas: identificación → RAG + resolución}.
  .4 ingestion/ \dtc{validación de audio/texto y normalización}.
  .3 platform/.
  .4 api/ \dtc{routers, autenticación, arranque}.
  .4 observability/ \dtc{trazas MLflow, latencia, presupuesto de tokens}.
  .4 worker/ \dtc{cola de tareas asíncrona en proceso}.
  .3 repositories/ \dtc{capa de acceso a datos que el pipeline consume}.
  .3 services/ \dtc{lógica de dominio: admin, graph, health, job}.
  .3 transcription/.
  .2 cli/.
  .2 core/.
  .3 config/.
  .3 models/ \dtc{nodos, relaciones, enumeraciones, modelos API}.
  .3 utils/ \dtc{retry, tokens, audio, logging}.
  .2 eval/ \dtc{tooling de calidad}.
  .3 calibration/ \dtc{meta-scorers para umbrales de confianza y retrieval}.
  .3 grounding/ \dtc{arnés del agente de verificación}.
  .3 grouping/ \dtc{arnés del agente de agrupación de decisiones}.
  .3 identification/ \dtc{arnés de extracción de nodos}.
  .3 resolution/ \dtc{arnés de resolución de relaciones}.
  .3 retrieval/ \dtc{arnés del motor de búsqueda}.
  .2 infra/.
  .3 blob\_store/.
  .3 knowledge\_store/ \dtc{nodos y relaciones}.
  .3 ops\_store/ \dtc{estado de trabajos y claves de API}.
  .3 secrets/ \dtc{resolución de credenciales}.
  .3 vector\_store/.
}
\caption{Estructura de paquetes de \textbf{Seshat}.}
\label{fig:package-tree}
\end{figure}

% ---------------------------------------------------------------------------
\section{Configuración centralizada}

Toda la configuración se centraliza en un único objeto \texttt{SeshatConfig}\footnote{La
referencia completa de parámetros y variables de entorno está disponible en
\url{https://github.com/josepmafe/seshat/blob/main/docs/configuration.md}.},
que hereda de \texttt{pydantic\_settings.BaseSettings}, con \texttt{\_\_} como separador de
anidamiento. Sus objetos anidados son subclases de \texttt{pydantic.BaseModel} planos de
modo que la resolución de variables de entorno ocurre exclusivamente en la raíz. El objeto
resultante es un \textit{singleton} inmutable cargado en el arranque del proceso.

Para los trabajos que requieren parámetros diferentes a los globales —por ejemplo, un umbral
de confianza más estricto o un modo automático activado—, la función \texttt{get\_request\_settings}
aplica una fusión profunda de un objeto \texttt{SeshatConfigOverride} opcional sobre el
\textit{singleton} base y devuelve un nuevo \texttt{SeshatConfig}. La fusión es recursiva a
cualquier profundidad: un campo no presente en el \textit{override} conserva su valor base en
todos los niveles. La configuración resultante —global o por trabajo— gobierna el comportamiento
de todas las etapas que se describen a continuación.

% ---------------------------------------------------------------------------
\section{Ingesta}
La ingesta comprende tres pasos: la validación del archivo de entrada, la transcripción de los
archivos de audio y la normalización de ambas rutas de entrada en un \texttt{TranscriptDocument}
común.

\subsection{Validación}
\label{sec:validacion-ingesta}

Los archivos de audio (.mp3, .wav, .m4a) pasan tres controles secuenciales antes de subirse a
blob storage:

\begin{enumerate}
  \item \textbf{Tamaño}: el archivo se rechaza si supera \texttt{max\_file\_bytes} (500~MiB por
        defecto).
  \item \textbf{Duración}: se extrae con \texttt{mutagen} y se rechaza si supera
        \texttt{max\_audio\_seconds} (dos horas por defecto).
  \item \textbf{Firma de bytes}: los primeros 12~bytes determinan el formato real
        (\textit{magic bytes}), que se contrasta con la extensión declarada; cualquier
        discrepancia se rechaza. La extensión por sí sola no es fiable al controlarla el cliente.
\end{enumerate}

El texto preformateado (.yaml, .yml, .json) solo requiere una comprobación de esquema y de
coherencia de la fecha de reunión, tras la cual el campo \texttt{content} se extrae directamente
como transcripción.

\subsection{Transcripción}

La transcripción de audio se delega en \texttt{AbstractTranscriber} —AssemblyAI por defecto—,
envuelta en \texttt{TrackingTranscriber} para acumular los segundos consumidos contra el
presupuesto de uso. El archivo original y la transcripción se persisten en blob storage como
artefactos de auditoría independientes de la base de conocimiento; si el trabajo falla más
adelante, ambos están disponibles para reintentar sin re-transcribir.

% ---------------------------------------------------------------------------
\section{Extracción multi-agente en dos pasadas}

La extracción es la fase de mayor coste computacional y el núcleo del sistema. Se articula en
dos pasadas secuenciales con un contrato estricto: todos los agentes de la primera pasada deben
completarse antes de que comience la segunda. Entre ambas se interpone la revisión humana opcional
(HITL): la primera pasada identifica y puntúa los nodos, y sólo para los nodos aprobados
—automática o manualmente— la segunda infiere las relaciones.

Todos los agentes comparten una base común que encapsula la invocación estructurada del LLM con
reintentos, y se descubren en tiempo de ejecución a partir de la configuración; habilitar o
deshabilitar un tipo de concepto es así un cambio de configuración y no de código. Sus prompts
de sistema, estáticos por tipo, se marcan además como cacheables a nivel de proveedor para
reducir coste y latencia. La arquitectura de clases de los agentes y el mecanismo de caché se
detallan en el anexo~\ref{anx:tecnico}.

Las cuatro subsecciones siguientes describen respectivamente la primera pasada —identificación y puntuación de
confianza—, y la segunda —recuperación y resolución—.

\subsection{Identificación}
\label{sec:agentes-identificacion}

La transcripción completa se envía en paralelo a cuatro agentes de identificación, uno por tipo
de concepto (figura~\ref{fig:identification-flow}). El \ac{MVP} no realiza segmentación previa: la
ventana de contexto de los modelos actuales es suficiente para transcripciones de reuniones
típicas. Cada agente recibe además un \textit{hint} ligero con los nodos de su mismo tipo ya
presentes en la base de conocimiento —título y resumen de los más recientes, acotado en
tokens—, lo que le permite ver ejemplos de nodos ya aprobados.

\begin{figure}[H]
  \centering
  \includegraphics[width=\textwidth]{identification_flow.drawio.png}
  \caption{Flujo de la primera pasada de extracción. La transcripción normalizada en la ingesta
           alimenta el \textit{fan-out} de agentes de identificación —el de decisiones incluye
           además el paso de agrupación—; sus salidas se deduplican, se puntúan por confianza
           —heurísticas y, opcionalmente, verificación por \textit{grounding}— y reciben un estado
           de aprobación (pasos descritos en la sección~\ref{sec:confidence-scoring}) antes de
           persistirse como nodos nuevos.}
  \label{fig:identification-flow}
\end{figure}

Cada agente devuelve una lista de nodos estructurados mediante salida forzada con esquemas
Pydantic. Cada nodo queda anclado a una cita literal del texto fuente que lo respalda, lo que
reduce la alucinación y hace auditable cada afirmación. El diseño de los prompts de
identificación combina varias técnicas complementarias que se detallan en el anexo~\ref{anx:tecnico}.

Para el tipo \texttt{decision} se introduce un paso adicional de agrupación: la lista plana de
conceptos se organiza en grupos temáticos, de modo que el \textit{pipeline} genere un único
nodo por decisión consolidada en lugar de múltiples nodos fragmentados. El paso es configurable
por tipo, pues sólo las decisiones —que tienden a dispersarse a lo largo de un intercambio— se
benefician de forma clara; ante un fallo del agrupador, cada concepto retrocede a un grupo
unitario y ningún nodo se pierde.

Tras la ejecución de los agentes, los nodos con título idéntico y mismo tipo se fusionan
(\textit{deduplicación within-meeting}): el nodo más reciente en la salida del agente —que se
asume representa el estado final de la deliberación— absorbe los anteriores.

Opcionalmente, cada agente de identificación puede envolverse en una segunda pasada de
auto-revisión que valida y filtra sus propias salidas antes de puntuarlas; el mecanismo se
describe en el anexo~\ref{anx:tecnico}.

\subsection{Puntuación de confianza y estado de aprobación}
\label{sec:confidence-scoring}

Cada nodo recibe, antes de pasar a revisión, una puntuación de confianza en $[0,1]$: una señal
continua, resultado de la media ponderada de tres métricas heurísticas. Cada métrica aplica
reglas léxico-sintácticas sobre el análisis morfológico, de dependencias y de entidades que
produce el modelo \texttt{en\_core\_web\_sm} de \texttt{spaCy}:

\begin{equation}
  \text{confianza} = 0{,}3 \cdot s_{\text{cita}} + 0{,}3 \cdot s_{\text{título}} + 0{,}4 \cdot s_{\text{descripción}}
\end{equation}

donde $s_{\text{cita}}$ es una señal que satura a las 35 palabras de cita fuente. $s_{\text{título}}$ mide
la especificidad del título combinando tres señales —presencia de entidades nombradas, calificadores y
longitud, en ese orden de peso— que reflejan la importancia relativa de estas señales para distinguir
enunciados específicos de generales~\cite{louis2011specificity}. $s_{\text{descripción}}$ penaliza
multiplicativamente el lenguaje de cobertura, la voz pasiva y el tiempo futuro; los dos primeros son
marcadores léxicos de incertidumbre ya empleados en la detección de \textit{hedging}~\cite{farkas2010conll}.

Opcionalmente, el \texttt{GroundingAgent} actúa como compuerta binaria independiente: un modelo
ligero de un proveedor distinto al de los agentes de identificación —la configuración lo exige
para evitar errores correlacionados— juzga si la descripción está respaldada por la cita fuente.
Un nodo que falla la verificación es rechazado o marcado como pendiente de revisión,
independientemente de su puntuación heurística.

La política de aprobación automática combina ambas señales: si la puntuación supera el umbral
configurado y la verificación pasa (o está desactivada), el nodo se aprueba automáticamente con
\texttt{approval\_method=THRESHOLD}; en caso contrario el nodo queda en estado \texttt{AWAITING\_REVIEW}.
Con \texttt{auto\_mode=True}, todos los nodos cuya puntuación supere el umbral establecido se
aprueban sin revisión humana; los que no son rechazados automaticamente.

\subsection{Recuperación aumentada}
\label{sec:rag-resolucion}

El módulo de recuperación (RAG) opera en la segunda pasada, una vez que los nodos de la primera
han sido consolidados y aprobados. A diferencia del uso habitual de RAG, orientado a
responder preguntas fundamentando la respuesta en los documentos previamente ingeridos, aquí la
recuperación localiza los nodos preexistentes de la base de conocimiento con los que cada nodo
nuevo podría relacionarse.

Para cada nodo nuevo, el \texttt{NodeRetriever} construye una consulta a partir de su título y
descripción, la envía al \texttt{SearchEngine} con el umbral de similitud configurado y expande
los resultados con los vecinos directos del grafo (profundidad~1); por último, un reranking
opcional reordena los candidatos por relevancia. Un presupuesto de contexto
acotado en tokens limita el número total de candidatos para no sobrepasar la ventana del modelo
de resolución. La figura~\ref{fig:retrieval-flow} resume el recorrido de un nodo fuente a través
del recuperador.

\begin{figure}[H]
  \centering
  \includegraphics[width=\textwidth]{retrieval_flow.drawio.png}
  \caption{Flujo de recuperación para un nodo fuente, en su configuración completa (modo
           \texttt{HYBRID}). El \texttt{SearchEngine} combina una pata semántica (búsqueda densa,
           con expansión \textit{multi-query} opcional) y una pata dispersa por palabras clave,
           fundidas mediante RRF. El \texttt{NodeRetriever} obtiene los nodos de la base de
           conocimiento, los expande con sus vecinos directos y opcionalmente aplica un reranking
           (recuadro discontinuo) antes de devolver los nodos objetivo.}
  \label{fig:retrieval-flow}
\end{figure}

Un nodo puede recuperarse por dos vías complementarias, cada una con sus propias limitaciones. La
búsqueda \textbf{densa} (\textit{dense}) compara la consulta con los nodos en el espacio de
\textit{embeddings} y recupera por proximidad semántica: reconoce paráfrasis y sinónimos aunque
no compartan vocabulario, pero tiende a difuminar los términos raros y específicos. La búsqueda
\textbf{dispersa} (\textit{sparse}) casa términos literales y sobresale justo donde la densa se
diluye —nombres propios, herramientas o identificadores técnicos que el \textit{embedding}
promedia—, aunque no capta reformulaciones. \texttt{SearchEngine} expone ambas, y su combinación,
en tres modos de búsqueda configurables:

\begin{itemize}
  \item \texttt{SEMANTIC}: búsqueda densa por similitud coseno sobre pgvector. Es el modo por
        defecto. Opcionalmente admite un \textit{fan-out} multi-consulta: un LLM genera variantes
        de la consulta original —más abstractas, más específicas o con distinto vocabulario— y sus
        resultados se funden con \ac{RRF}~\cite{cormack2009rrf}, ampliando la cobertura frente a una
        única formulación.
  \item \texttt{KEYWORD}: búsqueda dispersa. Un LLM ligero destila primero la consulta a 3--6
        palabras clave discriminantes y la búsqueda se realiza sobre ellas; sin este paso, la
        coincidencia literal —que exige la presencia conjunta de todos los términos— se ahogaría
        en el vocabulario genérico de una consulta larga y devolvería poco o nada.
  \item \texttt{HYBRID}: combina la pata densa —con su \textit{fan-out} multi-consulta si está
        activo— y la dispersa, fusionándolas mediante \ac{RRF}, que las agrega por la posición
        de cada nodo en cada lista y no por su puntuación. Esto es deliberado: las similitudes
        coseno y las de texto viven en escalas distintas y no son directamente comparables,
        mientras que el rango sí lo es.
\end{itemize}

Sobre la lista devuelta por cualquiera de los modos puede aplicarse un \textit{reranking}
opcional. Mientras que la búsqueda densa aproxima la relevancia por la distancia entre vectores
calculados de forma independiente para consulta y candidato, un modelo \textit{cross-encoder}
procesa ambos textos de forma conjunta y estima su relevancia directamente, lo que reordena los
candidatos con mayor precisión a costa de una llamada adicional a un proveedor externo.

Por último, la recuperación es intencionalmente \textbf{multi-tipo}: un nodo de tipo
\texttt{decision} puede necesitar resolver relaciones frente a nodos de tipo \texttt{risk} o
\texttt{open\_question}, de modo que restringir la búsqueda al mismo tipo privaría de candidatos
válidos a los agentes de resolución cruzada.

\subsection{Agentes de resolución}
\label{sec:agentes-resolucion}

La resolución de relaciones opera en dos familias paralelas: \textbf{mismo tipo} y
\textbf{tipo cruzado} (figura~\ref{fig:resolution-flow}). En ambas, las relaciones son
\textbf{dirigidas} y se infieren en un único sentido: desde cada nodo nuevo (fuente) hacia los
nodos preexistentes de la base de conocimiento (destino) recuperados en la fase anterior. Los
nodos de reuniones previas no se reevalúan entre sí; sólo la llegada de un nodo nuevo dispara la
inferencia de relaciones, lo que mantiene el coste acotado al conocimiento incremental de cada
trabajo.

%\vspace{1.5cm}

\begin{figure}[H]
  \centering
  \includegraphics[width=\textwidth]{resolution_flow.drawio.png}
  \caption{Flujo de la segunda pasada. Para cada nodo fuente de la primera pasada se recuperan
           sus nodos objetivo (recuperación descrita en la sección~\ref{sec:rag-resolucion}); el
           par fuente–objetivos se envía en paralelo a las familias de resolución de mismo tipo
           (un agente por tipo) y de tipo cruzado (hasta nueve pares activos). Un paso de
           validación mediante reglas heurísticas filtra las relaciones inferidas antes de
           persistirlas.}
  \label{fig:resolution-flow}
\end{figure}

Los agentes de mismo tipo clasifican cada nodo nuevo frente a los candidatos recuperados de la
base de conocimiento en una de las relaciones aplicables a su tipo o sin relación. Por ejemplo,
para un nodo de tipo \texttt{decision}, las relaciones permitidas del mismo tipo son
\texttt{supersedes}, \texttt{amends} y \texttt{conflicts\_with} (el conjunto completo por tipo se
recoge en la tabla~\ref{tab:rel-matrix} del anexo~\ref{anx:tecnico}). Además, el prompt establece
criterios operativos precisos: \texttt{supersedes} requiere que la decisión previa quede
funcionalmente obsoleta; \texttt{amends} implica que la decisión previa permanece vigente pero es
matizada o restringida; \texttt{conflicts\_with} requiere que ambas decisiones sean activas
simultáneamente e incompatibles. Ante ambigüedad entre \texttt{amends} y \texttt{supersedes}, el
sistema prefiere \texttt{amends} como clasificación menos destructiva.

Los agentes de tipo cruzado detectan relaciones entre nodos de tipos distintos:
\texttt{mitigates}, \texttt{blocks} y \texttt{resolves}. Por ejemplo, una \texttt{decision} puede
mitigar (\texttt{mitigates}) un \texttt{risk} o resolver (\texttt{resolves}) una
\texttt{open\_question}. El esquema de relaciones restringe qué pares de tipos pueden evaluarse:
de los 12 pares dirigidos posibles entre tipos distintos, el \ac{MVP} habilita los 9 cuyas fronteras
semánticas son más nítidas, materializados como instancias de agente fijas; los restantes quedan
fuera de alcance por ser más ambiguos de clasificar de forma fiable. La tabla~\ref{tab:rel-matrix}
del anexo~\ref{anx:tecnico} detalla los pares habilitados.

Cuando un agente tiene incertidumbre genuina entre dos relaciones, puede señalarlo mediante el
campo opcional \texttt{alt\_rel\_type}. Este campo es la interfaz entre el agente de resolución
y una segunda pasada de auto-revisión opcional que arbitra los casos ambiguos, análoga a la de
identificación y detallada igualmente en el anexo~\ref{anx:tecnico}.

Tras ambas llamadas paralelas, un paso de validación heurística descarta las relaciones
malformadas —dirección incorrecta, incompatibilidades de tipo, violaciones del esquema— y las
registra sin interrumpir el trabajo.

% ---------------------------------------------------------------------------
\section{Orquestación y ciclo de vida del pipeline}
\label{sec:orquestacion}

Las dos pasadas de extracción descritas hasta aquí no se invocan directamente: son etapas de un
\textit{pipeline} que la \ac{API} encola y un \textit{worker} asíncrono ejecuta. Cada \textit{pipeline}
avanza por la máquina de estados descrita en la figura~\ref{fig:job-lifecycle}, desde su recepción
hasta la persistencia de los resultados.

\subsection{API y creación del pipeline}

La \ac{API} REST (FastAPI) es el único punto de entrada programático al sistema. Un \textit{pipeline}
se crea enviando un archivo de entrada, que la \ac{API} valida y persiste junto con el registro del
trabajo antes de encolarlo para su procesamiento asíncrono: la petición retorna de inmediato, sin
esperar al resultado. El resto de operaciones del ciclo de vida —consultar estado, recuperar
resultados, aprobar la revisión o reintentar un \textit{pipeline} fallido— se exponen como
\textit{endpoints} sobre el mismo recurso; la referencia completa de la \ac{API} se recoge en el
anexo~\ref{anx:api-reference}. El \texttt{JobService} concentra esta lógica de dominio y se apoya
en repositorios que encapsulan el acceso a datos, de modo que ningún \textit{endpoint} consulta la
base de datos directamente.

En la creación se aplican además dos salvaguardas: un control de duplicados por
\textit{hash} del contenido —que evita re-procesar un archivo ya ingerido salvo forzado
explícito— y límites de tasa por usuario y de concurrencia global, que protegen el presupuesto
de tokens frente a envíos masivos.

\subsection{Worker asíncrono y máquina de estados}

El \textit{worker} consume los \textit{pipelines} encolados y los conduce a través de los estados
del ciclo de vida, actualizando el estado persistido en cada transición para que la \ac{API} pueda
reportarlo. La ejecución se divide en dos tramos separados por la revisión humana. El primero
—ingesta, transcripción e identificación— termina dejando el \textit{pipeline} en
\texttt{AWAITING\_REVIEW}. El segundo —recuperación, resolución y escritura— sólo arranca una vez
aprobados los nodos. Cuando el modo automático está activo, o cuando ningún nodo queda pendiente
de revisión, el segundo tramo se encadena sin intervención manual.

El \textit{worker} está diseñado para tolerar fallos sin dejar la base de conocimiento en un
estado inconsistente. La escritura final de nodos y relaciones es una única transacción, de modo
que un fallo en cualquier punto no deja estado parcial. Además, cuando un \textit{pipeline} falla
queda en estado \texttt{FAILED} conservando su archivo original; la operación de reintento lo
recupera del almacenamiento y vuelve a ejecutar el \textit{pipeline} completo desde la transcripción,
sin reutilizar el trabajo ya realizado —una reanudación consciente del progreso queda como trabajo
futuro (sección~\ref{sec:trabajo-futuro})—. Por último, los \textit{pipelines} interrumpidos por una
caída del proceso durante la escritura se detectan y marcan como recuperables en el siguiente arranque.

\subsection{Revisión humana (HITL)}

En \texttt{AWAITING\_REVIEW}, el \textit{pipeline} se detiene a la espera de que un revisor decida
el destino de los nodos extraídos. Las decisiones se envían en una única petición al
\textit{endpoint} de aprobación, que admite dos mecanismos combinables. El primero es una
aprobación masiva por umbral: aprueba de una vez todos los nodos pendientes cuya confianza supere
un valor dado. El segundo son decisiones individuales: aprobar, rechazar o editar un nodo concreto.
Cuando se usan juntos, la regla masiva se aplica primero y las decisiones individuales se superponen
a su resultado, de modo que el revisor puede, por ejemplo, aprobar en bloque todo lo que supere el
umbral y luego rechazar o corregir los casos puntuales que lo requieran. Las ediciones y aprobaciones
registran quién y cuándo, distinción que sostiene la trazabilidad descrita a continuación. Una vez
enviadas las decisiones, el \textit{pipeline} transita a la fase de resolución y el \textit{worker}
retoma el segundo tramo.

% ---------------------------------------------------------------------------
\section{Mecanismos de seguridad}

\subsection{Control de acceso y autenticación}

Todos los \textit{endpoints} de la \ac{API} requieren una clave de \ac{API} en la cabecera
\texttt{X-API-Key}. Las claves se almacenan como hashes bcrypt (factor de coste~12); el texto
plano nunca se persiste. La validación de cada petición usa comparación en tiempo constante para
prevenir ataques de temporización.

El modelo de roles es jerárquico; cada nivel hereda los permisos del anterior:

\begin{equation*}
  \texttt{viewer} \;\subset\; \texttt{reviewer} \;\subset\; \texttt{operator} \;\subset\; \texttt{admin}
\end{equation*}

\texttt{viewer} sólo consulta la base de conocimiento; \texttt{reviewer} es el nivel mínimo para
enviar un \textit{pipeline} y para decidir sobre los nodos en revisión —aprobarlos, rechazarlos o
editarlos—. Activar el modo automático se sitúa también en este nivel, pues equivale a aprobar en
bloque sin revisión manual. Las acciones más sensibles se reservan a niveles superiores: los
\textit{overrides} de configuración por trabajo requieren \texttt{operator}, mientras que las
operaciones de borrado y la re-ingesta forzosa —que elimina nodos existentes— requieren
\texttt{admin}. La provisión de claves de \ac{API} queda al margen de esta jerarquía: se realiza a
través de un router \texttt{/admin} separado, autenticado con una clave raíz (\textit{root})
almacenada en Secrets Manager en lugar de con una clave de usuario.

\subsection{Validación de entradas}

Más allá del control de acceso, el sistema trata como hostiles sus dos entradas no controladas. El
archivo del cliente se valida en la frontera de la \ac{API} por firma de bytes
(sección~\ref{sec:validacion-ingesta}) en fallo cerrado: una entrada inválida se rechaza antes de
crear el \textit{pipeline} y no deja ni trabajo encolado ni artefacto almacenado. El texto de la
transcripción, que llega a los agentes, se aísla frente a la inyección de prompts mediante los
delimitadores XML del anexo~\ref{anx:tecnico}, reforzados por la salida estructurada forzada:
aunque el modelo procese contenido malicioso, el resultado debe ajustarse al esquema Pydantic para
no ser descartado.

% ---------------------------------------------------------------------------
\section{Observabilidad y trazabilidad}
\label{sec:observabilidad}

La observabilidad es una preocupación transversal en un sistema construido sobre \ac{LLM}: las
llamadas a los modelos son costosas, no deterministas y ocurren en paralelo, por lo que sin
instrumentación es difícil diagnosticar un fallo, atribuir un coste o justificar una afirmación
extraída. \textbf{Seshat} aborda esto en tres planos: el trazado de la ejecución, la medición del
consumo con enforcement de presupuesto, y la trazabilidad del conocimiento producido.

\subsection{Trazado de la ejecución}

MLflow~3 actúa como \textit{backbone} de observabilidad. La instrumentación de las llamadas a los
agentes es automática: \texttt{mlflow.langchain.autolog()} captura prompts, respuestas, uso de
tokens y latencia sin código adicional en los agentes, emitiendo cada llamada como un \textit{span}
de \ac{OTel}, el estándar abierto de trazado distribuido. Cada \textit{pipeline} abre un
\textit{run} de MLflow al arrancar; su identificador se
expone en la consulta de estado del trabajo, de modo que el operador puede saltar directamente a la
traza completa de la ejecución.

\subsection{Medición del consumo y presupuesto}

Sobre el trazado automático, \textbf{Seshat} añade una capa propia de medición que acumula, por
etapa del \textit{pipeline}, los tokens de entrada, salida y caché del LLM, los tokens de
\textit{embeddings} y de \textit{reranking}, y los segundos de audio transcritos, junto con
percentiles de latencia por llamada (media, p95 y máximo). Estas métricas se registran en MLflow y
hacen comparables el coste y el rendimiento entre versiones del sistema.

Esta capa no es meramente informativa: aplica un presupuesto de consumo. Cada etapa acumula su uso
contra unos límites configurables y aborta la ejecución si los rebasa, lo que acota el gasto ante
transcripciones anómalas o bucles inesperados. El control de gasto real recae así sobre estos
límites de tokens, y no sobre tablas de precios que requerirían mantenimiento; los costes
monetarios se estiman de forma informativa en la interfaz de MLflow. Los detalles de
implementación de esta capa —su propagación transparente a los agentes concurrentes y los
envoltorios de medición— se describen en el anexo~\ref{anx:observabilidad}.

\subsection{Trazabilidad del conocimiento}

Al margen de la instrumentación de ejecución, cada nodo conserva la traza de su origen y de su
manipulación posterior. La cita literal del texto fuente (\texttt{quote\_anchors}) que originó el
nodo hace auditable cualquier afirmación extraída. Las correcciones humanas en revisión quedan
registradas en los campos \texttt{corrected\_by} y \texttt{corrected\_at}, distinguiendo el
contenido corregido del generado automáticamente, y las anulaciones (\textit{overrides}) de
operadores y administradores exigen justificación escrita, almacenada en
\texttt{correction\_reason}.

% ---------------------------------------------------------------------------
\section{Despliegue e infraestructura}

\subsection{Contenedores Docker}

El sistema se orquesta con Docker Compose en seis servicios: \texttt{postgres}
(PostgreSQL~17 con pgvector), \texttt{mlflow} (observabilidad, con \textit{backend} SQLite),
\texttt{localstack} (emulación de S3 y Secrets Manager), \texttt{migrate} (aplica las migraciones
de esquema mediante Alembic y termina), \texttt{seshat-api} (la \ac{API} FastAPI, que aloja además el \textit{worker}
asíncrono en proceso) y \texttt{seshat-ui} (la interfaz de revisión en Streamlit). El arranque
está ordenado por comprobaciones de salud: la \ac{API} sólo se inicia cuando Postgres, LocalStack y
MLflow están operativos y las migraciones han concluido.

La \ac{API} y las migraciones se construyen desde un mismo \textit{Dockerfile} multietapa —comparten la
capa de dependencias y difieren sólo en la etapa final—, de modo que ambas ejecutan exactamente el
mismo entorno. La configuración se inyecta por variables de entorno mediante el separador de
anidamiento de pydantic-settings; los secretos —incluida la cadena de conexión de Postgres, bajo
la clave \texttt{seshat/postgres\_url}— se resuelven en el arranque desde Secrets Manager, de modo
que el mismo código se ejecuta en desarrollo (LocalStack) y en producción (AWS real) sin ramas
condicionales.

\subsection{Integración continua}
\label{sec:ci}

Cada \textit{push} a la rama principal y cada \textit{pull request} dispara una canalización de
\ac{CI} (GitHub Actions) que actúa como guardián de calidad antes de la fusión. La canalización
arranca con una etapa de \textit{linting} —formato y reglas estáticas con \textit{ruff}, más
comprobación de tipos con \textit{mypy}— de la que dependen las demás: si el estilo o los tipos
fallan, no se malgasta cómputo en las pruebas. Superada esa puerta, las pruebas unitarias y las de
integración (sección~\ref{sec:tests}) se ejecutan en paralelo; estas últimas levantan los mismos
contenedores de Docker Compose de la sección anterior, de modo que se validan contra Postgres, LocalStack y
MLflow reales, y exigen una cobertura de ramas mínima del 80\,\%. Una última etapa, sólo en las \textit{pull requests}, publica
como comentario métricas de complejidad ciclomática y de mantenibilidad del código, dando
visibilidad a la deuda técnica en el momento de la revisión.

La canalización ejecuta deliberadamente todas las pruebas salvo las marcadas como dependientes de
un \ac{LLM} en vivo: reproducir la aleatoriedad y el coste de las llamadas reales al modelo en cada
\textit{commit} no es ni fiable ni económico. La evaluación de la calidad de los agentes recae por
ello en el \textit{gate} descrito en el capítulo~\ref{ch:evaluacion}, que se ejecuta de forma
deliberada fuera de la \ac{CI} y cuyo veredicto la \ac{API} comprueba en el arranque.
